\begin{center}
{\Large \bfseries \textcolor{lyred}{第五章 定积分}}
\end{center}

\vspace{6pt}
{\large \bfseries \textcolor{lyred}{第5.1 节 定积分的概念与性质}}

\vspace{4pt}
{\bfseries \textcolor{lyred}{一.定积分问题举例}}

{\bfseries \textcolor{lyred}{1.曲边梯形的面积}}

{\bfseries \textcolor{lyred}{2.变速直线运动的路程}}

{\bfseries \textcolor{lyred}{3.细棒的质量}}

\vspace{4pt}
{\bfseries \textcolor{lyred}{二.定积分的定义}}

设
()
f x 为[
]
,a b 上的有界函数,将[
]
,a b 任意分成n 段[
]
1,
i
i
x
x
-
,长度
i
i
i
x
x
x
\Delta 
-
=
-
, 
令
\{
\}
max
,
,
n
x
x
\lambda =
\Delta 
\Delta 
\Lambda 
,
[
]
1,
i
i
i
x
x
\psi 
-
\forall 
\in 
,作和
()
n
i
i
i
f
x
\psi 
=
\Delta 
\sum 
,若在无限细分[
]
,a b 的 
过程中,随着
\lambda \to 
,该和总是趋向于同一个常数I ,它只依赖于
()
f x 和[
]
,a b ,则 
称
()
f x 在[
]
,a b 上\textcolor{lyred}{可积},并记
()
b
a
I
f x dx
=\int 
,称为
()
f x 在[
]
,a b 上的\textcolor{lyred}{定积分}. 
()
n
i
i
i
f
x
\psi 
=
\Delta 
\sum 
为\textcolor{lyred}{积分和},
()
f x 为\textcolor{lyred}{被积函数},
()
f x dx 为\textcolor{lyred}{被积表达式}, x 为\textcolor{lyred}{积分变量}, 
a 为\textcolor{lyred}{积分下限},b 为\textcolor{lyred}{积分上限},[
]
,a b 为\textcolor{lyred}{积分区间}. 
\textcolor{lyred}{注}.定积分的值与积分变量的符号无关. 
\vspace{4pt}
{\bfseries \textcolor{lyred}{三.可积的条件}}

\textcolor{lyred}{定理}.[
]
,a b 上只有有限个间断点的有界函数必可积. 
\textcolor{lyred}{推论}.[
]
,a b 上的连续函数必可积. 
\vspace{4pt}
{\bfseries \textcolor{lyred}{四.定积分与不定积分的关系}}

\textcolor{lyred}{定理}.设
()
f x 在[
]
,a b 上连续,且
()
()
F
x
f x
'
=
,则
()
()
()
b
a
f x dx
F b
F a
=
-
\int 
. 
\vspace{4pt}
{\bfseries \textcolor{lyred}{五.定积分的几何与物理意义}}

{\bfseries \textcolor{lyred}{1.几何意义}}

(1)当
()
f x \ge 
时,
()
b
a
f x dx
A
=
\int 
;(2)当
()
f x \le 
时,
()
b
a
f x dx
A
=-
\int 
; 
(3)一般地,
()
b
a
f x dx
A
A
+
-
=
-
\int 
,即曲边梯形面积的代数和. 
\textcolor{lyred}{注}.当
()
f x 为奇函数时
()
a
a
f x dx
-
=
\int 
,为偶函数时
()
()
b
a
a
f x dx
f x dx
=
\int 
\int 
. 
\textcolor{lyred}{例}.
a
a
x dx
a
\pi 
-
=
\int 
,
(
)
x
x dx
x
dx
\pi 
-
=
-
-
=
\int 
\int 
. 
\textcolor{lyred}{例}.
tan
x
x dx
x
-
=
+
\int 
,
sin
x
xdx
-
=
\int 
. 
\textcolor{lyred}{例}.
[
]
1 cos
sin
2 sin
cos
xdx
x dx
xdx
x
\pi 
\pi 
\pi 
\pi 
\pi 
\pi 
-
-
-
=
=
=
-
=
\int 
\int 
\int 
. 
\textcolor{lyred}{例}. (
)
x
x
dx
x dx
x dx
-
-


=
=
=
=




\int 
\int 
\int 
. 
{\bfseries \textcolor{lyred}{2.物理意义}}

(1)若已知直线运动的速度
()
v
v t
=
,则位移
()
T
T
S
v t dt
=\int 
. 
(2)设细棒位于[
]
,a b 上,有连续的线密度
()
x
\rho 
,则质量
()
b
a
M
x dx
\rho 
=\int 
. 
\vspace{4pt}
{\bfseries \textcolor{lyred}{六.利用定积分计算数列极限}}

设
()
f x 在[
]
,a b 上连续,取
(
)
i
i
x
a
b
a
n
=
+
-
,则
()
()
lim
b
n
i
n
i
a
b
a
f x
f x dx
n
\to \infty 
=
-
=
\sum 
\int 
; 
特别地,
()
lim
lim
n
n
n
n
i
i
i
i
f
f
f x dx
n
n
n
n
\to \infty 
\to \infty 
=
=
-




=
=








\sum 
\sum 
\int 
. 
\textcolor{lyred}{例}.
lim
sin
sin
sin
lim
sin
sin
n
n
n
i
n
i
xdx
n
n
n
n
n
n
\pi 
\pi 
\pi 
\pi 
\pi 
\pi 
\to \infty 
\to \infty 
=


+
+
+
=
=
=




\sum 
\int 
\Lambda 
. 
\textcolor{lyred}{例}.设
(
)(
)
n
n
i
a
n
i
n
i
=
=
+-
+
\sum 
,求lim
n
n
a
\to \infty 
.\textcolor{lyred}{ }
解.
n
n
n
i
i
a
n
i
n
i
=
=
\le 
\le 
+
+-
\sum 
\sum 
,而
lim
lim
ln 2
n
n
n
n
i
i
dx
i
n
i
n
x
n
\to \infty 
\to \infty 
=
=
=
=
=
+
+
+
\sum 
\sum 
\int 
, 
ln 2
n
n
i
i
dx
i
n
i
n
x
n
=
=
=
=
=
-
+-
+
+
\sum 
\sum 
\int 
,故lim
ln 2
n
n
a
\to \infty 
=
. 
\vspace{4pt}
{\bfseries \textcolor{lyred}{七.定积分的性质}}

\textcolor{lyred}{规定}:
()
()
a
b
b
a
f x dx
f x dx
=-
\int 
\int 
,例如:
x dx
x dx
=-
=-
\int 
\int 
. 
\textcolor{lyred}{性质1(线性性)}.
()
()
()
()
b
b
b
a
a
a
f x
g x
dx
f x dx
g x dx
\alpha 
\beta 
\alpha 
\beta 
\pm 
=
\pm 




\int 
\int 
\int 
. 
\textcolor{lyred}{性质2(对积分区间的可加性)}.
()
()
()
b
c
b
a
a
c
f x dx
f x dx
f x dx
=
+
\int 
\int 
\int 
. 
\textcolor{lyred}{例}.
(
)
1 sin 2
cos
sin
cos
sin
xdx
x
x
dx
x
x dx
\pi 
\pi 
\pi 
-
=
-
=
-
=
\int 
\int 
\int 
(
)
(
)
cos
sin
sin
cos
2 2
x
x dx
x
x dx
\pi 
\pi 
\pi 
-
+
-
=
\int 
\int 
. 
\textcolor{lyred}{例}.
x dx
x dx
xdx
xdx
-
-
-
=
=-
+
=
\int 
\int 
\int 
\int 
. 
\textcolor{lyred}{性质3}.
(
)
b
a
k dx
k b
a
\cdots 
=
-
\int 
;特别地,
b
b
a
a
dx
dx
b
a
=
\cdots 
=
-
\int 
\int 
. 
\textcolor{lyred}{性质4}.设在[
]
,a b 上
()
f x \ge 
,则
()
b
a
f x dx \ge 
\int 
. 
\textcolor{lyred}{注}.若
()
[
]
,
f x
C a b
\in 
,
()
f x \ge 
,且不恒为0 ,则
()
b
a
f x dx >
\int 
. 
\textcolor{lyred}{推论}.设在[
]
,a b 上
()
()
f x
g x
\ge 
,则
()
()
b
b
a
a
f x dx
g x dx
\ge 
\int 
\int 
. 
\textcolor{lyred}{例}.证明:
lim
n
n
x
dx
x
\to \infty 
=
+
\int 
. 
证.由
n
n
x
dx
x dx
n
x
\le 
\le 
=
+
+
\int 
\int 
,即得,证毕. 
\textcolor{lyred}{推论}.
()
()
b
b
a
a
f x dx
f x dx
\le 
\int 
\int 
. 
\textcolor{lyred}{性质5(估值定理)}.设在[
]
,a b 上
()
m
f x
M
\le 
\le 
,则 
(
)
()
(
)
b
a
m b
a
f x dx
M b
a
-
\le 
\le 
-
\int 
. 
\textcolor{lyred}{性质6(积分中值定理)}.设
()
f x 在[
]
,a b 上连续,则存在
[
]
,a b
\psi \in 
,使得 
()
()(
)
b
a
f x dx
f
b
a
\psi 
=
-
\int 
,即
()
()
()
lim
b
n
i
n
i
a
f
f x dx
f x
b
a
n
\psi 
\to \infty 
=
=
=
-
\sum 
\int 
. 
\textcolor{lyred}{注}.该值称为
()
f x 在[
]
,a b 上的\textcolor{lyred}{平均值}. 
\textcolor{lyred}{积分第一中值定理}.设
()
f x , ()
g x 均在[
]
,a b 上连续,且()
g x \ge 
,则 
存在
[
]
,a b
\psi \in 
,使得
()()
()
()
b
b
a
a
f x g x dx
f
g x dx
\psi 
=
\int 
\int 
. 
\textcolor{lyred}{注}.“
[
]
,a b
\psi 
\exists \in 
”可以改进为“
(
)
,a b
\psi 
\exists \in 
”. 
\textcolor{lyred}{例}.设
()
f x 可导,
()
()
f
f x dx
=\int 
,证明:
(
)
0,1
\psi 
\exists \in 
,使得
()
f
\psi 
'
=
. 
证.
()
()
()
f
f
f
\eta 
\eta 
=
\cdots 
=
,
2 ,1
\eta 


\in 



,即得,证毕. 
\textcolor{lyred}{例}.设
()
f x 可导,
()
()
f
xf x dx
=\int 
,证明:存在
(
)
0,1
\psi \in 
,使得
()
()
f
f
\psi 
\psi 
\psi 
'
+
=
. 
证.
()
()
f
f
\eta 
\eta 
=
,
(
)
0,1
\eta \in 
,令
()
()
F x
xf x
=
,则
()
()
F
F \eta 
=
,证毕. 
\vspace{4pt}
{\bfseries \textcolor{lyred}{八.积分不等式}}

\textcolor{lyred}{例}.设
()
f x >
,连续,证明:
()
()
ln
ln
f x dx
f x dx
\le 
\int 
\int 
. 
证.
()
ln
lim
ln
limln
limln
n
n
n
n
n
n
n
i
i
i
i
i
i
f x dx
f
f
f
n
n
n
n
n
\to \infty 
\to \infty 
\to \infty 
=
=
=






=
=
\le 
=












\sum 
\sum 

\int 
()
ln lim
ln
n
n
i
i
f
f x dx
n
n
\to \infty 
=

=




\sum 
\int 
,证毕. 
\textcolor{lyred}{例}.设
()
sin
,
1,
x
x
f x
x
x

\neq 

=

=

,证明:
()
f x dx
\pi 
\pi 
\le 
\le 
\int 
. 
证.
()
f x 在0, 2
\pi 






上单调减少,故
()
f x
\pi \le 
\le ,即得,证毕. 
\textcolor{lyred}{例}.设
()
f x 可导,
()
f a =
,
()
f
x
M
'
\le 
,证明:
()
(
)
b
a
M
f x dx
b
a
\le 
-
\int 
. 
证.
()
()
()
()
(
)
b
b
b
b
a
a
a
a
f x dx
f x
f a dx
f
x
a dx
M
x
a dx
\psi 
'
=
-
=
-
\le 
-
\int 
\int 
\int 
\int 
,证毕. 
\textcolor{lyred}{例}.设
()
f x 连续,单调增加,0
\lambda 
<
<,证明:
()
()
f x dx
f x dx
\lambda 
\lambda 
<
\int 
\int 
. 
证.
()
()
(
)
()
()
f x dx
f x dx
f x dx
f x dx
\lambda 
\lambda 
\lambda 
\lambda 
\lambda 
\lambda 
-
=
-
-
=
\int 
\int 
\int 
\int 
(
)
()
(
)(
)
(
)
()
(
)
f
f
f
f
\lambda \lambda 
\psi 
\lambda 
\psi 
\lambda 
\lambda 
\lambda 
\psi 
\psi 
-
-
-
=
-
-
<




,证毕. 
\textcolor{lyred}{例(柯西积分不等式)}.设
()
f x 与()
g x 均在[
]
,a b 上连续,则 
()()
()
()
b
b
b
a
a
a
f x g x dx
f
x dx
g
x dx

\le 
\cdots 




\int 
\int 
\int 
. 
证.
(
)
,
t
\forall \in -\infty +\infty ,均有
()
()
b
a
f x
tg x
dx
-
\ge 




\int 
,于是由 
()
()()
()
b
b
b
a
a
a
t
g
x dx
t
f x g x dx
f
x dx
-
+
\ge 
\Rightarrow \Delta \le 
\int 
\int 
\int 
,即得,证毕. 
\textcolor{lyred}{注}.(1)
()
()
()
(
)
()
b
b
b
b
b
a
a
a
a
a
f x dx
f x dx
dx
f
x dx
b
a
f
x dx




=
\cdots 
\le 
\cdots 
=
-








\int 
\int 
\int 
\int 
\int 
. 
(1)设
()
f x >
,则
()
()
()
()
(
)
b
b
b
a
a
a
f x dx
dx
f x
dx
b
a
f x
f x




\cdots 
\ge 
\cdots 
=
-




\int 
\int 
\int 
. 
\textcolor{lyred}{补充练习 }
1.设
n
n
i
i
a
n
i
=
=
+
\sum 
,求lim
n
n
a
\to \infty 
.\textcolor{lyred}{ }
解.
(
)
lim
lim
ln 1
ln 2
n
n
n
n
i
i
x
n
a
dx
x
n
x
i
n
\to \infty 
\to \infty 
=


=
=
=
+
=


+


+



\sum 
\int 
. 
2.设
n
n
i
n
a
n
i
=
=
+
+
\sum 
,求lim
n
n
a
\to \infty 
. 
解.
(
)
n
n
n
i
i
n
n
a
n
i
n
i
=
=
\le 
\le 
+
+
+
\sum 
\sum 
,
(
)
lim
lim
n
n
n
n
i
i
n
n
n
i
i
n
\to \infty 
\to \infty 
=
=
=
+
+
+


+



\sum 
\sum 
, 
lim
lim
n
n
n
n
i
i
n
n
i
n
i
n
\to \infty 
\to \infty 
=
=
=
+


+



\sum 
\sum 
,故
lim
n
n
dx
a
x
\pi 
\to \infty 
=
=
+
\int 
. 
3.比较
sin
x
P
dx
x
\pi 
\pi 
-
=
+
\int 
,
sin
x
Q
dx
x
\pi 
\pi 
-
=
+
\int 
,
(
)
sin
cos
x
x
R
dx
x
\pi 
\pi 
-
-
=
+
\int 
的大小. 
解.
1 2sin cos
x
x
R
dx
dx
Q
x
x
\pi 
\pi 
\pi 
\pi 
-
-
-
=
=
>
>
+
+
\int 
\int 
,而
P =
,故R
Q
P
>
>
. 
4.设
sin
n
na
x
xdx
-
=\int 
,求lim
n
n
a
\to \infty 
. 
解.
n
n
na
x dx
x dx
n
-
\le 
=
=
+
\int 
\int 
,故lim
n
n
a
\to \infty 
=
. 
5.设
()
f x 连续,
()1
f
=,求
()
lim
n
n
n
n
n
n
f x dx
+
\to \infty 
-\int 
. 
解.原式
()
()
()
lim
3lim
n
n f
f
f
n
\psi 
\psi 
\psi 
\to \infty 
\to 
=
\cdots 
\cdots 
=
=
=
. 
6.设
()
f x 有连续的导数,求
(
)
(
)
lim
a
a
a
f x
a
f x
a
dx
a
\to 
-
+
-
-




\int 
. 
解.上式
(
)
(
)
()
()
2lim
4lim
a
a
f
a
f
a
f
f
a
\psi 
\psi 
\eta 
\to 
\to 
+
-
-
'
'
=
=
=
. 
7.设
()
f x 可导,
()
f x dx =
\int 
,证明:
(
)
0,1
\psi 
\exists \in 
,使得
()
()
f
f
\psi 
\psi 
\psi 
'
+
=
. 
证.存在
(
)
0,1
c\in 
,使得
()
f c =
;令
()
()
F x
xf x
=
,则
()
()
F
F c
=
=
,证毕. 
8.设
()
()
f x
x
x
f x dx
=
-
+\int 
,求
()
f x . 
解.设
()
f x dx
A
=
\int 
,则
()
f x
x
Ax
=
-
+
,故
(
)
A
x
Ax
dx
=
-
+
=
\int 
x dx
A x dx
A
A
\pi 
\pi 
-
+
=
+
\Rightarrow 
=
\int 
\int 
,于是
()
f x
x
x
\pi 
=
-
+
. 
9.设
()
()
()
f x
x
x
f x dx
f x dx
=
-
+
\int 
\int 
,求
()
f x . 
解.设
()
f x dx
A
=
\int 
,
()
f x dx
B
=
\int 
,则
()
f x
x
Ax
B
=
-
+
,于是 
(
)
A
x
Ax
B dx
A
B
=
-
+
=
-
+
\int 
,
(
)
B
x
Ax
B dx
A
B
=
-
+
=
-
+
\int 
,得 
A =
,
B =
,故
()
f x
x
x
=
-
+
. 
\vspace{6pt}
{\large \bfseries \textcolor{lyred}{第5.2 节 微积分基本公式}}

\vspace{4pt}
{\bfseries \textcolor{lyred}{一.积分上限函数及其导数}}

\textcolor{lyred}{定理}.设
()
f x 在[
]
,a b 上连续,则
()
()
x
a
x
f t dt
\Phi 
=\int 
在[
]
,a b 上可导,且 
()
()
x
f x
'
\Phi 
=
,
(
)
,
x
a b
\forall \in 
,而
()
()
a
f a
+'
\Phi 
=
,
()
()
b
f b
-'
\Phi 
=
. 
\textcolor{lyred}{注}.当
()
f x 在[
]
,a b 上可积时,
()
x
\Phi 
一定在[
]
,a b 上连续. 
\textcolor{lyred}{推论}.连续函数必有原函数. 
\textcolor{lyred}{推论}.
()
()
()
()
x
a
f t dt
f
x
x
\varphi 
\varphi 
\varphi 
'


'
=
\cdots 










\int 
,
()
()
()
()
b
x
f t dt
f
x
x
\zeta 
\zeta 
\zeta 
'


'
=-
\cdots 










\int 
;一般地, 
()
()
()
()
()
()
()
x
x
f t dt
f
x
x
f
x
x
\varphi 
\zeta 
\varphi 
\varphi 
\zeta 
\zeta 
'


'
'
=
\cdots 
-
\cdots 














\int 
. 
\textcolor{lyred}{例}.
sin
sin
cos
cos
cos
sin
x
t
x
x
x
d
e dt
x e
x e
dx

=
\cdots 
+
\cdots 




\int 
. 
\textcolor{lyred}{例}.
(
)
cos
cos
cos
2 sin
x
x
x
d
d
x
t
t dt
x
t dt
t
t dt
x
x
dx
dx




-
\cdots 
=
\cdots 
-
\cdots 
=








\int 
\int 
\int 
. 
\textcolor{lyred}{例}.设
()
,
1,
x
t
x
e dt
x
F x
x
x

\neq 

=

=

\int 
,证明:(1)
()
F x 连续;(2)
()
F x 可导. 
解.
()
lim
lim
lim
x
t
x
x
x
x
x
x
e dt
e
e
F x
x
\to 
\to 
\to 
-
=
=
=
\int 
, 
()
()
()
lim
lim
lim
x
x
t
t
x
x
x
x
x
e dt
e dt
x
F x
F
x
F
x
x
x
\to 
\to 
\to 
-
-
-
'
=
=
=
=
-
\int 
\int 
lim
lim
x
x
x
x
x
x
e
e
xe
xe
x
\to 
\to 
-
-
-
=
=
. 
\textcolor{lyred}{例}.设
()
f x 连续,
()
f x >
,证明:
()
()
()
x
x
tf t dt
F x
f t dt
=\int 
\int 
在(
)
0,\infty 上单增. 
证.
()
()
()
()
()
()
()(
)
()
()
x
x
x
x
x
xf x
f t dt
f x
tf t dt
f x
x
t f t dt
F
x
f t dt
f t dt
-
-
'
=
=
>












\int 
\int 
\int 
\int 
\int 
,证毕. 
\textcolor{lyred}{例}.设
()
f
x
'
\ge 
,令
()
()
x
a
F x
f t dt
x
a
=
-\int 
,证明:当x
a
\neq 
时,
()
F
x
'
\ge 
. 
证.
()
()(
)
()
(
)
()
()
(
)
x
x
a
a
f x
x
a
f t dt
f x
f t
dt
F
x
x
a
x
a
-
-
-




'
=
=
-
-
\int 
\int 
,即得,证毕. 
\textcolor{lyred}{注}.
()(
)
()
()
()(
)
x
a
f x
x
a
f t dt
f x
f
x
a
\psi 
-
-
=
-
-




\int 
,其中\psi 介于a 与x 之间. 
\vspace{4pt}
{\bfseries \textcolor{lyred}{二.牛顿-莱布尼茨公式}}

\textcolor{lyred}{定理}.设
()
f x 在[
]
,a b 上连续,且
()
()
F
x
f x
'
=
,则
()
()
()
b
a
f x dx
F b
F a
=
-
\int 
. 
\textcolor{lyred}{推论}.设
()
f x 在[
]
,a b 上连续,则存在
(
)
,a b
\psi \in 
,使得
()
()(
)
b
a
f x dx
f
b
a
\psi 
=
-
\int 
. 
\textcolor{lyred}{补充练习 }
1.设
()
1,
sin ,
x
x
f x
x
x

+
\le 
=
>

,
()
()
x
F x
f t dt
-
=\int 
,则在(
)
,
-\infty +\infty 上______. 
(A)
()
F x 是
()
f x 的原函数;  (B)
()
F x 不是
()
f x 的原函数,但可导; 
(C)
()
F x 不可导,但连续;    (D)
()
F x 不连续. 
解.
()
(
)
F
f
-
-'
=
=,
()
(
)
F
f
+
+'
=
=
,故选(C). 
2.设
()
(
)
sin
F x
x
t
t dt
x
-
=
-
-<
<
\int 
,求
()
F
x
''
. 
解.
()
sin
sin
sin
sin
x
x
x
x
F x
x
t dt
t
t dt
t
t dt
x
t dt
-
-
=
-
+
-
\int 
\int 
\int 
\int 
,故 
()
sin
sin
x
x
F
x
tdt
tdt
-
'
=
-
\int 
\int 
,
()
2sin
F
x
x
''
=
. 
3.将当
x
+
\to 
时的无穷小量:
cos
x
t dt
\alpha =\int 
,
tan
x
tdt
\beta =\int 
,
sin
x
t dt
\gamma =\int 
,按阶的 
由低到高的顺序排列. 
解.
cos
lim
lim
k
k
x
x
x
k
x
kx
\alpha 
+
+
-
\to 
\to 
=
\Rightarrow 
=,
2 tan
lim
lim
lim
k
k
k
x
x
x
x
x
x
k
x
kx
kx
\beta 
+
+
+
-
-
\to 
\to 
\to 
=
=
\Rightarrow 
=
, 
sin
lim
lim
lim
k
k
k
x
x
x
x
x
x
k
x
kx
kx
\gamma 
+
+
+
-
-
\to 
\to 
\to 
=
=
\Rightarrow 
=
,故排序是
, ,
\alpha \gamma \beta . 
4.设
()
f
x
'
连续,
()
f
=
,令
()
()
,
0,
x
tf t dt
x
F x
x
x

\neq 

=

=

\int 
,讨论
()
F
x
'
的连续性. 
解.当
x \neq 
时,
()
()
()
()
()
x
x
xf x
x
x
tf t dt
x f x
tf t dt
F
x
x
x
\cdots 
-
\cdots 
-
'
=
=
\int 
\int 
,而 
()
()
()
()
()
()
()
lim
lim
lim
lim
x
x
x
x
x
tf t dt
F x
F
xf x
f x
f
F
x
x
x
x
\to 
\to 
\to 
\to 
'
-
'
=
=
=
=
=
\int 
; 
()
()
()
()
()
()
lim
lim
lim
x
x
x
xf x
x f
x
xf x
f
x
f
F
x
x
\to 
\to 
\to 
'
'
'
+
-
'
=
=
=
,故连续. 
5.设
()
[
]
,
f
x
C a b
\in 
,单调减少,证明:
()
()
b
b
a
a
a
b
xf
x dx
f
x dx
+
\le 
\int 
\int 
. 
证.令
()
()
(
)
()
x
x
a
a
F x
tf t dt
a
x
f t dt
=
-
+
\int 
\int 
,
()
(
)
()
()
x
a
F
x
x
a f x
f t dt
'
=
-
-
=
\int 
()
()
x
a
f x
f t
dt
-
\le 




\int 
,而
()
F a =
,证毕. 
6.设
()
f x 连续,
()
f x dx =
\int 
,证明:存在
(
)
0,1
\psi \in 
,使得
()
()
f x dx
f
\psi 
\lambda 
\psi 
+
=
\int 
. 
证.令
()
()
x
x
F x
e
f t dt
\lambda 
=
\int 
,则
()
()
F
F
=
=
,即得,证毕. 
7.设
()
f x 与()
g x 均在[
]
,a b 上连续,证明:存在
(
)
,a b
\psi \in 
,使得 
()
()
()
()
b
a
f
g x dx
g
f x dx
\psi 
\psi 
\psi 
\psi 
=
\int 
\int 
. 
证.令
()
()
()
x
b
a
x
F x
f t dt
g t dt
=
\cdots 
\int 
\int 
,则
()
()
F a
F b
=
=
,即得,证毕. 
\vspace{6pt}
{\large \bfseries \textcolor{lyred}{第5.3 节 定积分的换元法和分部积分法}}

\vspace{4pt}
{\bfseries \textcolor{lyred}{一.定积分的换元法}}

\textcolor{lyred}{例}.
(
)
(
)
x
x dx
x
d x
t dt
\pi 
-
=
-
-
-
=
-
=
\int 
\int 
\int 
. 
\textcolor{lyred}{例}.
sin
sin
sin
cos
sin
cos
sin
cos
x
xdx
x
x dx
x
xdx
x
xdx
\pi 
\pi 
\pi 
\pi 
\pi 
-
=
=
-
=
\int 
\int 
\int 
\int 
sin
sin
sin
sin
xd
x
xd
x
u du
u du
\pi 
\pi 
\pi 
-
=
-
=
\int 
\int 
\int 
\int 
. 
\textcolor{lyred}{例}.
(
)
2sin
4sin
1 cos2
2sin
sin
2cos
x
t
x
t
t
dx
d
t
tdt
dt
t
x
\pi 
\pi 
\pi 
\pi 
=
-
=
=
=
=
-
-
\int 
\int 
\int 
\int 
. 
\textcolor{lyred}{例}.
x u
dx
u
du
u
x
x
u
=


=-
=-
+
=
-


+
+
\int 
\int 
. 
\textcolor{lyred}{例}.
(
)
t
x
x
t
t
dx
d
t
dt
t
x
=
+


+
+
-
=
=
+
=


+


\int 
\int 
\int 
. 
\textcolor{lyred}{例}.
(
)
(
)
t
x
x
dx
t d t
t
t
dt
=
+
+
=
\cdots 
-
=
-
=
-
\int 
\int 
\int 
. 
\textcolor{lyred}{例}.求
(
)
sin
x
n
d
x
t dt
dx
-
\int 
. 
解.
(
)
(
)
sin
sin
sin
x
x
u x t
n
n
n
x
x t dt
ud x u
udu
=-
-
=
-
=
\int 
\int 
\int 
,故
(
)
sin
sin
x
n
n
d
x
t dt
x
dx
-
=
\int 
. 
\textcolor{lyred}{例}.设
()
f x 连续,且
(
)
1 cos
x
tf x
t dt
x
-
=-
\int 
,求
()
f x . 
解.
(
)
(
)
()
()
()
1 cos
x
x
x
u x t
x
tf x
t dt
x
u f u du
x
f u du
uf u du
x
=-
-
=-
-
=
-
=-
\int 
\int 
\int 
\int 
,故 
()
()
()
()
()
sin
sin
cos
x
x
f u du
xf x
xf x
x
f u du
x
f x
x
+
-
=
\Rightarrow 
=
\Rightarrow 
=
\int 
\int 
. 
\textcolor{lyred}{例}.设
()
f x 连续,单调增加,0
\lambda 
<
<,证明:
()
()
f x dx
f x dx
\lambda 
\lambda 
<
\int 
\int 
. 
\textcolor{lyred}{证}.
()
(
)
()
f x dx
f
t dt
f t dt
\lambda 
\lambda 
\lambda 
\lambda 
=
<
\int 
\int 
\int 
,证毕. 
\vspace{4pt}
{\bfseries \textcolor{lyred}{二.定积分的分部积分法}}

\textcolor{lyred}{定理}.设()
u x 与()
v x 均在[
]
,a b 上有连续的导数,则 
()()
()()
()
()
b
b
b
a
a
a
u x v x dx
u x v x
v x u
x dx
'
'
=
-




\int 
\int 
,即
[
]
b
b
b
a
a
a
udv
uv
vdu
=
-
\int 
\int 
. 
\textcolor{lyred}{例}.
[
]0
cos
sin
sin
sin
x
xdx
xd
x
x
x
xdx
\pi 
\pi 
\pi 
\pi 
=
=
-
=-
\int 
\int 
\int 
. 
\textcolor{lyred}{例}.
(
)
1 sin
sin
cos
1 sin
cos
cos
cos
cos
x
x
xdx
x
x
xd
x
dx
dx
xd
x
x
x
x
x
\pi 
\pi 
\pi 
\pi 
\pi 
\pi 
\pi 
-
-
-
-
=
=-
=
=
=
+
\int 
\int 
\int 
\int 
\int 
(
)
2ln
2 \pi 
-
+
+
. 
\textcolor{lyred}{例}.
ln
ln
ln
4ln 4
4ln 4
x dx
xd
x
x
x
dx
x
x
x




=
=
-
=
-
=
-




\int 
\int 
\int 
. 
\textcolor{lyred}{例}.
[
]
arcsin
arcsin
2 6
x
xdx
x
x
dx
x
x
\pi 
\pi 


=
-
=
\cdots 
+
-
=
+
-


-
\int 
\int 
. 
\textcolor{lyred}{例}.
(
)
ln
sin ln
sin
sin
sin1
cos
sin1
cos
e
t
x
t
t
t
t
x dx
e
tdt
tde
e
e
tdt
e
tde
=
=
=
=
-
=
-
=
\int 
\int 
\int 
\int 
\int 
(
)
sin1
cos1 1
sin
sin1
cos1 1
t
e
e
e
tdt
e
e
-
+-
=
-
+
\int 
. 
\textcolor{lyred}{例}.设
()
x
t
f x
e
dt
-
=\int 
,求
()
I
f x dx
=\int 
. 
解.
()
()
1 1
x
x
I
xf x
xf
x dx
xe
dx
e
e
-
-




'
=
-
=-
=
=
-












\int 
\int 
. 
\textcolor{lyred}{例}.设
()
sin
x
t
f x
dt
t
\pi 
=
-
\int 
,求
()
I
f x dx
\pi 
=\int 
. 
解.
()
()
sin
sin
sin
t
x
I
xf x
xf
x dx
dt
x
dx
xdx
t
x
\pi 
\pi 
\pi 
\pi 
\pi 
\pi \pi 
\pi 
'
=
-
=
-
=
=




-
-
\int 
\int 
\int 
\int 
. 
\textcolor{lyred}{例}.求
sin
cos
n
n
nI
xdx
xdx
\pi 
\pi 
=
=
\int 
\int 
. 
证.
(
)
sin
cos
sin
cos
cos
1 sin
cos
n
n
n
nI
xd
x
x
x
x
n
x
xdx
\pi 
\pi 
\pi 
-
-
-


=-
=-
+
\cdots 
-
=


\int 
\int 
(
)
(
)
(
)
(
)
sin
1 sin
n
n
n
n
n
n
n
x
x dx
n
I
n
I
I
I
n
\pi 
-
-
-
-
-
-
=
-
-
-
\Rightarrow 
=
\int 
,故 
(
)
(
)
1 !!
,
!!
sin
cos
1 !!,
!!
n
n
n
n
n
k
n
I
xdx
xdx
n
n
k
n
\pi 
\pi 
\pi 
-

\cdots 
=

=
=
=
-

=
+

\int 
\int 
. 
\vspace{4pt}
{\bfseries \textcolor{lyred}{三.定积分的特殊技巧}}

\textcolor{lyred}{定理}.设
()
f x 在[
]
,a a
-
上连续,则(1)
(
)
()
()
a
a
f
x
f x
f x dx
-
-
=-
\Rightarrow 
=
\int 
; 
(2)
(
)
()
()
()
a
a
a
f
x
f x
f x dx
f x dx
-
-
=
\Rightarrow 
=
\int 
\int 
. 
\textcolor{lyred}{例}.
(
)
ln
1 cos
sin
2 sin
x
x
x dx
x dx
xdx
\pi 
\pi 
\pi 
\pi 
\pi 
-
-


+
+
+
-
=
=
=




\int 
\int 
\int 
. 
\textcolor{lyred}{例}.
(
)
sin
x
x
x
dx
dx
x
dx
x
x
\pi 
-
+
=
=
-
-
=
-
+
-
+
-
\int 
\int 
\int 
. 
\textcolor{lyred}{例}.
(
)
(
)
x
x
x dx
x
x
dx
u
u du
\pi 
-
-
=
-
-
=
+
-
=
\int 
\int 
\int 
. 
\textcolor{lyred}{例}.设
()
f x 连续,则
()
(
)
b
b
a
a
f x dx
f a
b
x dx
=
+
-
\int 
\int 
. 
证.
()
(
)(
)
(
)
b
a
b
x a b t
a
b
a
f x dx
f a
b
t d a
b
t
f a
b
t dt
=+-
=
+
-
+
-
=
+
-
\int 
\int 
\int 
,证毕. 
\textcolor{lyred}{例}.
(
)
(
)
sin
cos
f
x dx
f
x dx
\pi 
\pi 
=
\int 
\int 
. 
\textcolor{lyred}{例}.
cos
sin
tan
cos
sin
cos
sin
n
n
n
n
n
n
n
dx
x
x
dx
dx
x
x
x
x
x
\pi 
\pi 
\pi 
\pi 
=
=
=
+
+
+
\int 
\int 
\int 
. 
\textcolor{lyred}{例}.
cos
cos
cos
cos
x
x
u
u
x
u
e
e
e
e
dx
du
\pi 
\pi 
\pi 
-
-
=-
-
-
=
=
\int 
\int 
. 
\textcolor{lyred}{例}.
(
)
tan
ln 1
tan
ln 1
tan
ln 1
1 1 tan
u
x dx
u
du
du
u
\pi 
\pi 
\pi 
\pi 


-




+
=
+
-
=
+
=






+\cdots 






\int 
\int 
\int 
(
)
ln 2
ln 2
ln
ln 1
tan
tan
du
x dx
u
\pi 
\pi 
\pi 
\pi 


=
-
+
=


+


\int 
\int 
. 
\textcolor{lyred}{例}.
x
x
x
u
u
u
u
u
u
u
u
x
xdx
u
du
udu
du
de
du
e
e
e
e
e
e
e
e
e
e
e
e
-
-
-
-
-
-
=
=
-
=
=
=
+
+
+
+
+
+
\int 
\int 
\int 
\int 
\int 
\int 
arctan
arctan
arctan
te
e
e
e
e
e




=
-








. 
\textcolor{lyred}{例}.设
()
f x 连续,则
(
)
(
)
sin
sin
xf
x dx
f
x dx
\pi 
\pi 
\pi 
=
\int 
\int 
. 
证.
(
)
(
)
(
)(
)
(
)
(
)
sin
sin
sin
x
t
xf
x dx
t f
t d
t
t f
t dt
\pi 
\pi 
\pi 
\pi 
\pi 
\pi 
\pi 
=-
=
-
-
=
-
\int 
\int 
\int 
,证毕. 
\textcolor{lyred}{例}.
sin
sin
sin
1 cos
1 cos
1 cos
x
x
x
x
x
du
dx
dx
dx
x
x
x
u
\pi 
\pi 
\pi 
\pi 
\pi 
\pi 
\pi 
-
-
=
=
=
=
+
+
+
+
\int 
\int 
\int 
\int 
. 
\textcolor{lyred}{例}.设
()
f x 连续,则
()
()
(
)
a
a
a
f x dx
f x
f
x
dx
-
=
+
-




\int 
\int 
. 
\textcolor{lyred}{例}.
sin
sin
sin
sin
x
x
x
x
x
x
dx
dx
xdx
e
e
e
\pi 
\pi 
\pi 
\pi 
\pi 
-
-


=
+
=
=


+
+
+


\int 
\int 
\int 
. 
\textcolor{lyred}{例}.设
()
f x 连续,则
()
()
(
)
a
a
f x dx
f x
f a
x
dx
=
+
-




\int 
\int 
. 
证.
()
()
()
()
(
)(
)
a
a
a
a
a
a
f x dx
f x dx
f x dx
f x dx
f a
t d a
t
=
+
=
+
-
-
=
\int 
\int 
\int 
\int 
\int 
()
(
)
a
a
f x dx
f a
t dt
+
-
\int 
\int 
,证毕. 
\textcolor{lyred}{例}.
(
)
(
)
sin
sin
f
x dx
f
x dx
\pi 
\pi 
=
\int 
\int 
. 
\textcolor{lyred}{例}.
sin
sin
sin
n
n
n
x
xdx
xdx
xdx
\pi 
\pi 
\pi 
\pi 
\pi 
=
=
\int 
\int 
\int 
. 
{\bfseries \textcolor{lyred}{3.周期性}}

\textcolor{lyred}{定理}.设
()
(
)
,
f x
C
\in 
-\infty +\infty ,有周期T ,则
()
()
a T
T
a
f x dx
f x dx
+
=
\int 
\int 
. 
\textcolor{lyred}{例}.证明:
sin
sin
x
e
xdx
\pi 
>
\int 
. 
证.
(
)
(
)
sin
sin
sin
sin
sin
sin
sin
sin
x
x
x
x
e
xdx
e
xdx
e
x
e
x
dx
\pi 
\pi 
\pi 
\pi 
-
-


=
=
+
-
=


\int 
\int 
\int 
(
)
sin
sin
sin
x
x
e
e
xdx
\pi 
-
-
>
\int 
,证毕. 
\textcolor{lyred}{例}.设
()
f x 为具有周期T 的连续奇函数,证明:
()
x
f t dt
\int 
也有周期T . 
证.
()
()
()
()
()
T
x T
x
x T
x
x
T
f t dt
f t dt
f t dt
f t dx
f t dt
+
+
-
=
+
=
+
\int 
\int 
\int 
\int 
\int 
,证毕. 
\textcolor{lyred}{注}.设
()
f x 具有周期T ,则
()
()
a nT
T
a
f x dx
n
f x dx
+
=
\int 
\int 
. 
\textcolor{lyred}{例}.
(
)
1 sin 2
cos
sin
cos
sin
20 2
xdx
x
x dx
x
x dx
\pi 
\pi 
\pi 
\pi 
\pi 
-
-
+
=
+
=
+
=
\int 
\int 
\int 
. 
\textcolor{lyred}{例}.
(
)
1 sin 2
sin
cos
sin
cos
20 2
xdx
x
x dx
x
x dx
\pi 
\pi 
\pi 
\pi 
\pi 
-
=
-
=
-
=
\int 
\int 
\int 
. 
\textcolor{lyred}{例}.
(
)
sin
sin
sin
sin
x
t
x
x dx
t
t dt
x
x dx
x dx
\pi 
\pi 
\pi 
\pi 
\pi 
\pi 
\pi 
=
-
=
-
\Rightarrow 
=
=
\int 
\int 
\int 
\int 
sin
sin
x dx
xdx
\pi 
\pi 
\pi 
\pi 
\pi 
=
=
\int 
\int 
. 
\textcolor{lyred}{补充练习 }
1.设
()
,
1 cos
,
x
x
x
f x
xe
x
-

-\le 
<
+
=

\ge 

,求
(
)
I
f x
dx
=
-
\int 
. 
解.
()
cos
tan
1 cos
t
t
dt
t
t
I
f t dt
te
dt
d
e
dt
t
e
-
-
-
-
-
=
=
+
=
+
=
-
+
+
\int 
\int 
\int 
\int 
\int 
. 
2.设
()
f x 连续,
()
lim
x
f x
x
\to 
=,令
()
(
)
x
F x
tf x
t
dt
=
-
\int 
,求
()
lim
x
F x
I
x
\to 
=
. 
解.
()
()
x
u x
t
F x
f u du
=
-
=
\int 
,故
()
(
)
lim
lim
x
x
xf x
F
x
I
x
x
\to 
\to 
'
=
=
=
. 
3.设
()
f x 连续,
()
lim
x
f x
x
\to 
=,
()
()
F x
f tx dt
=\int 
,讨论
()
F
x
'
的连续性. 
解.
()
()
x
u tx
f u du
F x
x
=
=\int 
,
()
()
()
()
lim
lim
1 lim 2
x
x
x
x
xf x
f u du
f x
F
x
x
x
\to 
\to 
\to 
-
'
=
=-
=
\int 
, 
()
()
()
()
lim
lim
lim 2
x
x
x
x
f u du
F x
f x
F
x
x
x
\to 
\to 
\to 
'
=
=
=
=
\int 
,故连续. 
4.设
()
f x 连续,
()
f
\neq 
,求
(
)
()
(
)
lim
x
x
x
x
t f t dt
I
x
f x
t dt
\to 
-
=
-
\int 
\int 
. 
解.
()
()
()
()
()
()
()
()
()
lim
lim
lim
x
x
x
x
x
x
x
x
x
x
x f t dt
tf t dt
f t dt
f t dt
x
I
x f u du
f u du
xf x
f u du
f x
x
\to 
\to 
\to 
-
=
=
=
=
+
+
\int 
\int 
\int 
\int 
\int 
\int 
\int 
()
()
()
()
()
()
lim
x
f
f
f
f x
f
f
\psi 
\psi 
\to 
=
=
+
+
. 
5.设
()
f x 连续,证明:
()
(
)
()
x
t
x
f u du dt
x
t f t dt


=
-




\int \int 
\int 
. 
证一.左式
()
()
()
x
x
t
t
x
f u du dt
t
f u du
tf t dt




=
=
-
=








\int \int 
\int 
\int 
右式. 
证二.当
x =
时成立; 左式求导
()
x
f u du
=\int 
, 
右式求导
()
()
()
()
x
x
f t dt
xf x
xf x
f t dt
=
+
-
=
\int 
\int 
,即得,证毕. 
6.
(
)(
)
(
)
(
)
x
x
x
x
x
x
x
x
e
e
dx
x e
e
dx
xd e
e
e
-
-
-
-
+
-
=
-
=
+
=
\int 
\int 
\int 
. 
7.
cos
sin
cos
sin
sin
cos
x
x
x
x
x
x
e
e
dx
dx
e
e
e
e
\pi 
\pi 
\pi 
=
=
+
+
\int 
\int 
. 
8.求
cos
cos
cos
cos
cos
sin
x
x
xdx
x
xdx
x
xdx
\pi 
\pi 
\pi 
\pi 
\pi 
\pi 
-
=
-
=
=
\int 
\int 
\int 
. 
9.
(
)
(
)
sin 2
tan
sin 2
tan
sin
6 sin
x
x
dx
x
x
dx
udu
udu
\pi 
\pi 
\pi 
\pi 
\pi 
\pi 
\pi 
+
+
-
-
+
=
+
=
=
=
\int 
\int 
\int 
\int 
3 1
6 4 2 2
\pi 
\pi 
\cdots 
\cdots 
\cdots 
=
. 
\vspace{6pt}
{\large \bfseries \textcolor{lyred}{第5.4 节 反常积分}}

\vspace{4pt}
{\bfseries \textcolor{lyred}{一.无穷限的反常积分}}

\textcolor{lyred}{例}.考虑
x
e dx
+\infty 
-
\int 
,它的几何意义是曲线
x
y
e-
=
与x 正半轴间区域的面积A ,而 
(
)
lim
lim 1
t
x
t
t
t
A
e dx
e
-
-
\to +\infty 
\to +\infty 
=
=
-
=
\int 
,故
x
e dx
+\infty 
-
=
\int 
. 
\textcolor{lyred}{定义}.若
()
lim
t
t
a
f x dx
\to +\infty \int 
存在,则称\textcolor{lyred}{反常积分}
()
a
f x dx
+\infty 
\int 
\textcolor{lyred}{收敛},其积分值 
()
()
lim
t
t
a
a
f x dx
f x dx
+\infty 
\to +\infty 
=
\int 
\int 
;反之,称
()
a
f x dx
+\infty 
\int 
\textcolor{lyred}{发散}; 
类似地,
()
()
lim
b
b
t
t
f x dx
f x dx
\to -\infty 
-\infty 
=
\int 
\int 
,
()
()
()
f x dx
f x dx
f x dx
+\infty 
+\infty 
-\infty 
-\infty 
=
+
\int 
\int 
\int 
. 
\textcolor{lyred}{注}.设
()
f x 连续,并且
()
()
F
x
f x
'
=
,则
()
()
a
a
f x dx
F x
+\infty 
+\infty 
=



\int 
, 
()
()
b
b
f x dx
F x
-\infty 
-\infty 
=



\int 
,
()
()
f x dx
F x
+\infty 
+\infty 
-\infty 
-\infty 
=



\int 
.(假定极限均存在) 
\textcolor{lyred}{例(p-积分)}.讨论
p dx
x
+\infty 
\int 
的敛散性. 
解.当
p =时,
lim
lim ln
t
t
t
dx
dx
t
x
x
+\infty 
\to +\infty 
\to +\infty 
=
=
=+\infty 
\int 
\int 
,发散; 
当
p <时,
lim
lim 1
t
p
p
p
t
t
dx
dx
t
x
x
p
p
+\infty 
-
\to +\infty 
\to +\infty 
=
=
-
=+\infty 
-
-
\int 
\int 
,发散; 
当
p >时,
lim
lim 1
t
p
p
t
t
dx
dx
t
x
x
p
p
p
+\infty 
-
\to +\infty 
\to +\infty 
=
=
-
=
-
-
-
\int 
\int 
,收敛. 
\textcolor{lyred}{注}.对
ln2
ln p
p
dx
du
x
x
u
+\infty 
+\infty 
=
\int 
\int 
,有相同的结果. 
\textcolor{lyred}{例}.
(
)
arctan
3 3
dx
dx
x
x
x
x
\pi 
\pi 
\pi 
+\infty 
+\infty 
+\infty 
+


=
=
=
-
=


+
+
+
+


\int 
\int 
. 
\textcolor{lyred}{例}.
xdx
xdx
xdx
du
du
x
x
x
u
u
+\infty 
+\infty 
+\infty 
-\infty 
-\infty 
+\infty 
=
+
=
+
+
+
+
\int 
\int 
\int 
\int 
\int 
,发散. 
\textcolor{lyred}{例}.
(
)
ln
ln 2
dx
x
dx
x
x
x
x
x
+\infty 
+\infty 
+\infty 



=
-
=
=




+
+
+




\int 
\int 
. 
\textcolor{lyred}{例}.
(
)
sec
sec
cos
1 sin
sin
sec
tan
x
t
dx
d
t
tdt
t d
t
t
t
x
x
\pi 
\pi 
\pi 
\pi 
\pi 
\pi 
+\infty 
=
=
=
=
-
=
-
-
\int 
\int 
\int 
\int 
. 
\textcolor{lyred}{例}.
x t
dx
t dt
x dx
d x
x
t
x
x
x
x
=
+\infty 
+\infty 
+\infty 
+\infty 
+


=
=
=
-
=


+
+
+




-
+




\int 
\int 
\int 
\int 
1 arctan
2 2
du
u
u
\pi 
+\infty 
+\infty 
-\infty 
-\infty 


=
=


+


\int 
. 
\textcolor{lyred}{例}.
x
x
x
x
xe
dx
xde
xe
e
dx
+\infty 
+\infty 
+\infty 
+\infty 
-
-
-
-


=-
=-
+
=


\int 
\int 
\int 
. 
\textcolor{lyred}{例}.
ln
ln
ln
x
x
dx
xd
dx
x
x
x
x
x
+\infty 
+\infty 
+\infty 
+\infty 
+\infty 






=
-
=-
+
=
+-
=












\int 
\int 
\int 
. 
\vspace{4pt}
{\bfseries \textcolor{lyred}{二.无界函数的反常积分}}

\textcolor{lyred}{例}.考虑
1 dx
x
\int 
,它的几何意义是
y
x
=
,
y =
,
x =
,
x =围成区域的面积,而 
(
)
lim
lim 2
t
t
t
A
dx
t
x
+
+
\to 
\to 
=
=
-
=
\int 
,故
dx
x
=
\int 
. 
\textcolor{lyred}{定义}.若在
()
U c
\forall 
内
()
f x 均无界,则称c 为\textcolor{lyred}{无界间断点},或者\textcolor{lyred}{瑕点}. 
\textcolor{lyred}{定义}.设
()
(
]
,
f x
C a b
\in 
, x
a
=
为其无界间断点(这无所谓),若
()
lim
b
t
a
t
f x dx
+
\to \int 
存在, 
则称\textcolor{lyred}{反常积分}
()
b
a
f x dx
\int 
\textcolor{lyred}{收敛},其积分值
()
()
lim
b
b
t
a
a
t
f x dx
f x dx
+
\to 
=
\int 
\int 
; 
类似地,若x
b
=
为
()
f x 的无界间断点,则
()
()
lim
b
t
t
b
a
a
f x dx
f x dx
-
\to 
=
\int 
\int 
; 
若
(
)
,
c
a b
\in 
为
()
f x 的无界间断点,则
()
()
()
b
c
b
a
a
c
f x dx
f x dx
f x dx
=
+
\int 
\int 
\int 
. 
\textcolor{lyred}{注}.设c 为
()
f x 的无界间断点,
()
()
F
x
f x
'
=
,则
()
()
c
c
a
a
f x dx
F x
-
=



\int 
, 
()
()
b
b
c
c
f x dx
F x
+
=



\int 
,
()
()
()
b
b
c
a
c
a
f x dx
F x
F x
-
+
=
+








\int 
.(假定极限均存在) 
\textcolor{lyred}{例(p-积分)}.讨论
p dx
x\int 
的敛散性. 
解.当
p =时,
lim
lim ln
t
t
t
dx
dx
t
x
x
+
+
\to 
\to 
=
=
-
=+\infty 
\int 
\int 
,发散; 
当
p >时,
lim
lim
p
p
p
t
t
t
dx
dx
t
x
x
p
p
+
+
-
\to 
\to 
=
=
-
=+\infty 
-
-
\int 
\int 
,发散; 
当
p <时,
lim
lim
p
p
p
t
t
t
dx
dx
t
x
x
p
p
p
+
+
-
\to 
\to 
=
=
-
=
-
-
-
\int 
\int 
,收敛. 
\textcolor{lyred}{例}.
(
)
2 arcsin
dx
d
x
x
x
x
x
\pi 
-
+


=
=
=


-
-
\int 
\int 
. 
\textcolor{lyred}{例}.
(
)
(
)
(
)
tan
2 cos
sec
dx
d
x
du
d
t
tdt
t
x x
x
u
\pi 
\pi 
+\infty 
+\infty 
+\infty 
=
=
=
=
=
+
+
+
\int 
\int 
\int 
\int 
\int 
. 
\textcolor{lyred}{例}.
1 arctan
u
t
dx
dx
dt
du
u
u
t t
x x
x
x
\pi 
+\infty 
+\infty 
+\infty 
+\infty 
+\infty 
=
-


=
=
=
=
=


+
-


-
-
\int 
\int 
\int 
\int 
. 
\textcolor{lyred}{例}.
dx
dx
dx
dx
dx
x
x
x
x
x
x
x
x
=
+
=
+
=
-
-
-




-
-
-
-








\int 
\int 
\int 
\int 
\int 
(
)
(
)
arcsin 2
ln
ln 2
x
x
x
x
\pi 
-
+




-
+
-
+
-
=
+
+












. 
\textcolor{lyred}{例}.
(
)
lim
lim
t
t
dx
dx
dx
x
x
x
x
x
-
+
\to 
\to 
-
-




=
+
=
-
-+-
-
-
=+\infty 








\int 
\int 
\int 
,发散. 
\textcolor{lyred}{例}.
[
]
ln
ln
ln
lim
ln
x
xdx
x
x
xd
x
x
x
+
+
\to 
=
-
=
-
-=-
\int 
\int 
. 
\textcolor{lyred}{例}.
(
)
(
)
ln 1
ln 1
x
x
dx
x
x
dx
x
-


-
=
-
+
=


-


\int 
\int 
(
)
(
)
(
)
ln 1
ln
ln 4
lim
1 ln 1
ln 4
x
x
x
x
x
x
x
x
-
-
\to 
+


-
+
-
=
+
-
-
-
=
-


-


. 
\textcolor{lyred}{例}.
ln
ln
ln
ln
x t
x
t
x
t
dx
d
dt
dx
x
t
t
x
t
=
+\infty 
+\infty 
+\infty 
+\infty 
=
=
=-
+
+
+
+
\int 
\int 
\int 
\int 
,故
ln
x dx
x
+\infty 
=
+
\int 
. 
\textcolor{lyred}{注}.这个解法有问题,没有讨论收敛性. 
\textcolor{lyred}{补充练习 }
1.设
()
(
)
,
,
x
x
f x
x
x
x

\le 
+

=

>
+

,求
()
()
x
F x
f t dt
-\infty 
=\int 
. 
解.当
x \le 
时,
()
arctan
x
dt
F x
x
t
\pi 
-\infty 
=
=
+
+
\int 
; 
当
x >
时,
()
(
)
(
)
2arctan
x
x
dt
dt
d t
F x
x
t
t
t
t
\pi 
\pi 
-\infty 
=
+
=
+
=
+
+
+
+
\int 
\int 
\int 
. 
2.
(
)
ln
1 ln 2
dx
x
dx
x
x
x
x
x
x
x
+\infty 
+\infty 
+\infty 
+




=
-
+
+
=
-
=-




+
+




\int 
\int 
. 
3.
(
)
(
)
sin
1 sin
x
x
x
x
dx
dx
t
t dt
x
x
\pi 
\pi 
+
+
=
=
+
=+
-
-
\int 
\int 
\int 
. 
4.
(
)
(
)
(
)
sin
1 sin
π
u
t
u
x
x
dx
dx
du
t
dt
x
x
u
x
\pi 
\pi 
=
-
-
+
=
=
=
+
=
-
-
-
-
\int 
\int 
\int 
\int 
. 
5.
(
)(
)(
)
(
)(
)
x t
dx
t dt
dx
x
x
x
t
t
\alpha 
\alpha 
\alpha 
\pi 
\alpha 
=
+\infty 
+\infty 
+\infty 
\ge 
=
=
=
+
+
+
+
+
\int 
\int 
\int 
. 
6.已知
sin
x dx
x
\pi 
+\infty 
=
\int 
,求(1)
sin cos
x
x dx
x
+\infty 
\int 
;(2)
sin x dx
x
+\infty 
\int 
. 
解.(1)
sin cos
sin 2
sin
x
x
x
u
dx
d x
du
x
x
u
\pi 
+\infty 
+\infty 
+\infty 
=
=
=
\int 
\int 
\int 
; 
(2)
sin
sin
2sin
cos
sin
x
x
x
x
dx
xd
dx
x
x
x
x
\pi 
+
+\infty 
+\infty 
+\infty 
+\infty 


-
-
=
=
+
=




\int 
\int 
\int 
. 